\section{Introduction}
\label{sec:introduction}

Integrated sensing and communication (ISAC) allows a wireless network to reuse
spectrum, hardware, and signal-processing resources for transmission and
environmental sensing~\cite{liu_dual_functional_6g,gonzalez_revolution_2024}.
UAV-enabled ISAC further exploits mobility and favorable line-of-sight geometry
to extend sensing coverage~\cite{mu_uav_isac,meng_uav_isac_challenges}. When
multiple UAVs observe the same target from different bistatic geometries, their
soft decisions can provide complementary evidence~\cite{meng_cooperative_isac_2025}.

This cooperation is not free. Each observation is produced by a sensing chain
$i\!\rightarrow\!q\!\rightarrow\!j$ and must then be delivered from receiver
$j$ to a fusion UAV. A strong echo may contribute little if its reporting link
is unreliable, its delay--Doppler (DD) energy is poorly captured, or its packet
cost exceeds its marginal detection value~\cite{yuan_dd_isac,
keskin_mimo_otfs_isac,chen_node_selection_backhaul_2025}. Existing
node-selection and backhaul-aware studies address parts of this coupling, while
distributed-detection models often abstract away the bistatic geometry and
OTFS DD structure~\cite{chawla_massive_mimo_dd,
chawla_parameter_detection_mimo_wsn,meng_network_level_isac_2024}.

A further architectural choice has received less attention: \emph{where should
reports about each target meet?} A target-agnostic rule that maximizes a UAV's
aggregate incoming rate tends to assign every target to the same central
receiver. That concentration changes all reporting rates, finite-blocklength
success probabilities, and even the feasible candidate sets seen by the inner
selector. Thus, improving only the link-ranking formula cannot correct an
unfavorable fusion destination.

We study prior-assisted confirmation or re-detection and introduce target-local
fusion planning. For every target, the fusion UAV is chosen using the predicted
target position rather than the unknown truth. The resulting target-specific
reporting graph is passed to an adaptive two-stage C2F selector whose final
greedy replay is aligned with the implemented detector. This separation is
deliberate: the fusion rule changes the information available to the selector,
whereas the fair marginal-gain objective and resource constraints remain fixed.

Our main contributions are as follows.

\begin{itemize}
	\item We formulate joint target-specific fusion planning and triplet-level
	soft-report selection for multi-UAV bistatic OTFS-ISAC. The fusion destination
	depends only on the predicted target belief, preventing truth leakage into the
	scheduler.

	\item We combine this outer planning layer with communication-error calibration,
	an adaptive coarse frontier, waveform-derived DD refinement, and detector-aligned
	greedy replay. The inner fair utility and communication price are unchanged,
	which isolates the architectural role of fusion placement.

	\item We provide a reproducible V1 evaluation: a 1000-trial paired main
	comparison, a controlled fusion-rule ablation, reporting-structure diagnostics,
	and sweeps over belief error, rate threshold, deployment area, communication
	price, blocklength, and fleet size.

	\item We identify the operating boundary rather than reporting a single favorable
	point. The results distinguish measured serial latency from a prospective
	parallel-scheduling lower bound and distinguish confirmatory MC=1000 evidence
	from exploratory MC=100 sweeps.
\end{itemize}

The remainder of the paper presents the system model in
Section~\ref{sec:system_model}, the target-local C2F method in
Section~\ref{sec:proposed_algorithm}, the V1 evidence in
Section~\ref{sec:simulation_results}, and conclusions in
Section~\ref{sec:conclusion}.
